\begin{landscape}
\footnotesize

\begin{longtable}{p{6cm} p{6cm} p{4cm} p{3cm} p{6cm}}

\caption{Research Questions} \\

\toprule
\textbf{RQ} & \textbf{Aims and scope} & \textbf{Methodological components} & \textbf{Expected outputs} & \textbf{Grounding references} \\
\midrule
\endfirsthead

\multicolumn{5}{l}{\tablename\ \thetable\ -- Continued} \\
\toprule
\textbf{RQ} & \textbf{Aims and scope} & \textbf{Methodological components} & \textbf{Expected outputs} & \textbf{Grounding references} \\
\midrule
\endhead

\midrule
\multicolumn{5}{r}{Continued on next page} \\
\endfoot

\bottomrule
\endlastfoot

RQ1. How can a PDF-to-structured-ESG transformation pipeline convert Indonesian/English sustainability reports and filings into a governance- and regulation-aligned, sentence- and paragraph-level ECS (ESG Content Structure) representation that supports ABSA at the sentence level? & design and evaluate an end-to-end pipeline that (i) performs robust PDF parsing, (ii) extracts and normalizes textual content, tables, and figures, (iii) aligns content to a joint ESG taxonomy anchored in GRI/SASB/TCFD concepts, and (iv) outputs a structured JSON/NER-ready representation with sentence-level section identifiers and regulatory tag mappings. & document layout analysis, OCR quality handling, cross-language alignment, topic-taxonomy mapping, and sentence-level evidence tagging. & a reproducible, auditable data representation that enables sentence-level ABSA and regulator-oriented tracing of evidence. & Climate-focused ABSA pipelines and document-to-structure transformations emphasize the need to anchor text to regulatory concepts and to manage multi-language inputs in structured formats \cite{10.21203/rs.3.rs-3600821/v1,10.2139/ssrn.3796152,10.1111/1911-3846.13069,10.1111/1911-3846.12832,10.1002/bse.4089,10.1371/journal.pone.0288052,10.1051/e3sconf/202343602004}. Foundational work on ClimateBERT and domain-adapted transformers informs the feasibility of accurate sentence-level extraction within climate and ESG contexts \cite{10.21203/rs.3.rs-3600821/v1,10.48550/arxiv.2510.01222,10.1111/1475-679x.12625,10.18653/v1/2023.emnlp-main.975} \\
\addlinespace
RQ2. How should ESG be categorized by aspect/pillar, sentiment, and tone in bilingual Indonesian/English corporate disclosures to enable fine-grained ABSA, while preserving cross-language comparability and regulatory alignment? & define and validate a multi-dimensional annotation schema that (i) maps sentences to ESG aspects (environment, social, governance), (ii) assigns aspect-level sentiment polarity and evaluative tone, and (iii) accounts for document-level tone as a contextual backdrop. The schema should support cross-language consistency and interoperability with GRI/SASB/TCFD mappings. & ontology design, cross-language term mapping (Indonesian-English), annotation guidelines, inter-annotator agreement assessment, and a pilot corpus annotated in both languages. & a reference corpus and a formal taxonomy enabling comparability across Indonesian and English disclosures and enabling downstream model training for ABSA. & ABSA literature stresses sentence- and aspect-level extraction and cross-lingual considerations, with cross-lingual ABSA surveys highlighting the need for consistent aspect taxonomies across languages \cite{10.1016/j.inffus.2025.103073,10.1111/eufm.12509,10.48550/arxiv.2511.00024,10.3390/math12193119,10.1002/csr.70089}. FinBERT/ClimateBERT-style domain models support sentence-level interpretation in ESG domains \cite{10.1111/1911-3846.12832,10.21203/rs.3.rs-3600821/v1,10.48550/arxiv.2510.01222}. \\
\addlinespace
RQ3. Do tone-based ABSA outputs differ meaningfully from ClimateBERT-style label classifications for climate-related disclosures, and what is the relationship between detected tone and climate-specific targets or commitments? & compare traditional tone classification (document/section-level polarity) with climate-focused label outputs produced by ClimateBERT-like models, including assessments of alignment with explicit climate targets (e.g., net-zero, emissions reductions). Investigate whether tone signals correlate with or diverge from target-specific classifications. & parallel labeling experiments using (i) a tone-focused ABSA pipe (global polarity) and (ii) a climate-target-oriented classifier (ClimateBERT-NetZero-like or ClimateBERT variants), followed by statistical analyses of concordance, discordance, and predictive validity for reported targets. & empirical criteria and guidelines for choosing or combining tone and climate-focused labels to maximize decision-usefulness and regulatory alignment. & ClimateBERT and zero-shot climate classification literature demonstrate how domain-specific transformers capture climate-related signals beyond generic sentiment, and several studies compare climate-target detection with broader tone signals \cite{10.48550/arxiv.2510.01222,10.2139/ssrn.3796152,10.1371/journal.pone.0288052,10.1051/e3sconf/202343602004,10.1108/dta-01-2024-0065}. Foundational climate-disclosure detection work shows that target specificity requires specialized models for robust performance \cite{10.21203/rs.3.rs-3600821/v1,10.1111/1475-679x.12625,10.18653/v1/2023.emnlp-main.975}. \\
\addlinespace
RQ4. What weaknesses arise in ABSA extraction outputs, and how can a diagnostics framework detect and quantify extraction errors (e.g., misidentified aspects, incorrect sentiment polarity, or missed regulatory mappings) to inform model improvement? & develop a sentence-level ABSA diagnostics framework that (i) identifies common error modes (noise, misclassification, lexical drift, cross-language misalignment), (ii) quantifies error types with interpretable metrics, and (iii) provides actionable feedback loops to improve model components (lexicon updates, annotation guidelines, cross-language alignment). & error taxonomy development, curated test sets with known ground truth, per-sentence calibration analyses, ablation studies, and model-retraining protocols. & a robust quality assurance regime for ESG ABSA systems, including dashboards for developers and regulators to monitor performance and track improvements over time. & ABSA and ESG ML literature repeatedly highlight challenges in explainability, calibration, and domain adaptation; cross-lingual ABSA studies stress the need for rigorous evaluation across languages to ensure reliable outputs \cite{10.1002/csr.70089,10.2139/ssrn.3738629,10.1108/mf-01-2023-0020,10.1016/j.inffus.2025.103073,10.1111/eufm.12509,10.48550/arxiv.2511.00024,10.3390/math12193119}. Explainability-focused integration with ABSA is recommended to support regulator-facing outputs \cite{10.1002/csr.70089,10.2139/ssrn.3738629,10.1108/mf-01-2023-0020}. & \\
\addlinespace
RQ5. How can documentation and visualization practices be designed to maximize reproducibility and auditability of ESG ABSA experiments, including annotation schemas, model configurations, prompts/templates, and evaluation results? & propose comprehensive documentation standards and visualization schemes that (i) capture data provenance (source documents, language, version), (ii) describe annotation schemas and inter-annotator agreement metrics, (iii) enumerate model architectures, hyperparameters, datasets, and prompts/templates, (iv) present evaluation metrics and error analyses, and (v) provide reproducible notebooks and data-sharing guidelines. & a reproducibility blueprint, standardized metadata schemas, version-controlled annotation guidelines, and visualization dashboards for performance, bias, and calibration across languages and frameworks. & a reproducibility package including annotated corpora, code, prompts/templates, and evaluation reports, enabling independent replication and regulatory audit. & across climate and ESG NLP, systematic documentation and reproducibility considerations are emphasized; works on climate claim detection and zero-shot analyses illustrate the need for transparent methodology and datasets \cite{10.2139/ssrn.3796152,10.1371/journal.pone.0288052,10.1051/e3sconf/202343602004}. The broader ESG literature also stresses standardized reporting of framework mappings and cognitive traceability \cite{10.1108/medar-11-2020-1097,10.1108/aaaj-06-2021-5330,10.1108/aaaj-02-2020-4445} \\
\addlinespace
RQ6. What is the stability of ABSA outputs across cross-model and cross-prompt configurations, and what ensemble or verification strategies yield the most reliable and regulator-friendly results in Indonesian/English bilingual ESG analysis? & quantify output stability across multiple models (e.g., generic ABSA, FinBERT-based ESG classifiers, ClimateBERT variants, multilingual transformers) and across prompts/templates, and evaluate ensemble strategies (e.g., majority voting, self-consistency checks, retrieval-augmented approaches) for reducing variance and increasing reliability. & prompt engineering experiments, cross-model comparisons under controlled data splits, stability metrics (e.g., per-sentence label agreement, variance in aspect extraction, consistency of regulatory mappings), and verifier modules that cross-check outputs against regulatory lexicons and numerical indicators. & practical guidance for practitioners on selecting robust configurations, with quantified trade-offs between performance, stability, and interpretability. &  \\
\addlinespace


\end{longtable}

\end{landscape}